% ============================================================
% 《素数计数与黎曼猜想》参考习题 · 第二部（第 7–9 章）
% 编译：XeLaTeX，连续两次。书内编号由同目录 书稿标签.aux 提供；
% 书稿改版后运行 更新书稿标签.sh 再编译，编号自动随书稿更新。
% ============================================================
\documentclass[11pt,a4paper]{article}
\usepackage[left=2.6cm,right=2.6cm,top=2.6cm,bottom=2.6cm]{geometry}
\usepackage{fontspec}
\usepackage{xeCJK}
\setCJKmainfont[ItalicFont=Songti SC]{Songti SC}
\setCJKsansfont[ItalicFont=Songti SC]{Songti SC}
\setCJKmonofont[ItalicFont=Songti SC]{Songti SC}
\usepackage{amsmath,amssymb,mathtools}
\usepackage{booktabs,array,enumitem,fancyhdr,xcolor}
\usepackage{listings}
\usepackage{xr-hyper}
\usepackage[hidelinks]{hyperref}
\externaldocument[B-]{书稿标签}

% ---------- 记号（与书稿一致） ----------
\DeclareMathOperator{\Li}{Li}
\DeclareMathOperator{\li}{li}
\DeclareMathOperator{\Res}{Res}
\newcommand{\RR}{\mathbb{R}}
\newcommand{\CC}{\mathbb{C}}
\newcommand{\dd}{\,\mathrm{d}}
\renewcommand{\Re}{\operatorname{Re}}
\renewcommand{\Im}{\operatorname{Im}}

% ---------- 书内引用：编号取自书稿，不生成外部链接 ----------
\newcommand{\bref}[1]{\ref*{B-#1}}

% ---------- 难度标记 ----------
\newcommand{\lvI}{$\bigstar$\,基础}
\newcommand{\lvII}{$\bigstar\bigstar$\,推导}
\newcommand{\lvIII}{$\bigstar\bigstar\bigstar$\,延伸}

% ---------- 题目、提示、解答的统一格式 ----------
\newcommand{\Q}[3]{%
  \par\addvspace{1.1em}\noindent
  \hypertarget{q:#1}{}\textbf{习题 #1}\quad\textbf{#3}\hfill{\small #2}\par
  \nopagebreak\smallskip}
\newcommand{\Hint}[1]{%
  \par\addvspace{0.8em}\noindent
  \hypertarget{h:#1}{}\textbf{#1}\quad{\small（\hyperlink{q:#1}{题目}\ $\cdot$\ \hyperlink{s:#1}{解答}）}\quad}
\newcommand{\Sol}[2]{%
  \par\addvspace{1.2em}\noindent
  \hypertarget{s:#1}{}\textbf{习题 #1\quad #2}\hfill{\small（\hyperlink{q:#1}{题目}）}\par
  \nopagebreak\smallskip}
\newcommand{\toHint}[1]{\hfill{\small\hyperlink{h:#1}{提示}\ $\cdot$\ \hyperlink{s:#1}{解答}}}

\setlist{itemsep=2pt,topsep=3pt,leftmargin=2.2em}
\setlist[enumerate,1]{label=(\alph*)}
\setlength{\parindent}{2em}
\setlength{\parskip}{0.25em}
\linespread{1.18}
\allowdisplaybreaks

\lstset{basicstyle=\ttfamily\small,columns=fullflexible,frame=single,
  framerule=0.3pt,rulecolor=\color{gray!60},xleftmargin=1em,xrightmargin=1em,
  breaklines=true,showstringspaces=false}

\newcommand{\versiondate}{2026 年 10 月 23 日（修订）}
\pagestyle{fancy}
\fancyhf{}
\fancyhead[L]{\small《素数计数与黎曼猜想》参考习题 · 第二部}
\fancyhead[R]{\small 初稿 \versiondate}
\fancyfoot[C]{\small\thepage}
\renewcommand{\headrulewidth}{0.4pt}

\renewcommand{\contentsname}{目录}

\begin{document}

\begin{center}
{\LARGE\bfseries 参考习题 · 第二部\par}
\vspace{0.6em}
{\large $\zeta$ 的解析延拓与函数方程（第 7–9 章）\par}
\vspace{0.6em}
{\normalsize 李成桐《素数计数与黎曼猜想：$\zeta$ 函数、零点与显式公式》配套资源\par}
\vspace{0.3em}
{\small 版本：\versiondate（初稿，尚未逐题复核）\par}
\end{center}

\vspace{0.8em}
\section*{使用说明}

本书是导读型专著，正文论证自足完整，不设习题。这里的题目供希望通过练习检查和加深理解的读者自选使用，不做也不影响阅读正文。

\begin{itemize}
  \item \textbf{三档难度。}
    \lvI：检查定义、结论与具体计算；
    \lvII：独立完成正文中的关键推导或其变式；
    \lvIII：需要少量正文以外的论证或数值实验，供有余力的读者选做。
  \item \textbf{题目、提示、解答分开。}
    建议先独立思考，再看提示，最后对照解答。三部分之间可点击题号互相跳转。
  \item \textbf{书内编号。}
    题中「定理~\bref{thm:reflection}」「式~(\bref{eq:asymmetric_fe_alt})」「第~\bref{sec:riemann_first_method}~节」等编号均自动取自当前书稿；书稿改版后重新编译本文件，编号随之更新。
  \item \textbf{固定题号。}
    每题编号（如 9.3）一经发布不再改动；删去的题号留空，不由新题顶替，新题接续编号。反馈时请注明题号。
  \item \textbf{与第一部的联系。}
    个别题目用到第一部习题的结论（如由第一部习题 4.2 的解析表示展开得到的 $\zeta(1+u)=1/u+\gamma_E+O(u)$），题中会注明。
  \item \textbf{主选与可选。}
    星级表示完成难度，不等于题目的作用。初读时建议先做主选题；习题 7.2、7.3（完整估计部分）、9.3、9.6 为可选延伸，供希望深入的读者选做。
  \item \textbf{数值。}
    数值题的参考值用 Python 的 \texttt{mpmath} 计算（常规高精度计算，不是区间认证）；「$\approx$」表示四舍五入。
\end{itemize}

本次修订（2026 年 10 月 23 日）依据读者审查意见改正了习题 7.1 的题设与 8.2 的论证，补充了 7.4、8.1、9.1、9.5 等题的关键理由，统一 Euler 常数记号为 $\gamma_E$，并新增习题 9.7（Riemann 第二法）。

本批题目为初稿，题目可解性与解答仍待作者逐题复核。发现错误或表述不清，请在配套仓库的 Issues 中注明题号与具体位置。

\tableofcontents

% ============================================================
\newpage
\part*{一、题目}
\addcontentsline{toc}{part}{一、题目}
% ============================================================

\section*{第 \bref{ch:gamma_function} 章\quad $\Gamma$ 函数及其定义域的扩展}
\addcontentsline{toc}{section}{\texorpdfstring{第 \protect\bref{ch:gamma_function} 章\quad }{}\texorpdfstring{$\Gamma$}{Γ} 函数及其定义域的扩展}

\Q{7.1}{\lvI}{负半整数处的值与留数}
\begin{enumerate}
  \item 由递推公式~(\bref{eq:gamma-recurrence}) 与推论~\bref{cor:gamma-half}，求 $\Gamma(\tfrac52)$、$\Gamma(-\tfrac12)$、$\Gamma(-\tfrac32)$。
  \item 用余元公式~(\bref{eq:reflection-formula}) 核对 $\Gamma(-\tfrac12)\Gamma(\tfrac32)$。
  \item 写出 $\Gamma$ 在 $s=-1$、$s=-2$ 处的留数，并说明：在相邻两个极点所分隔的每个实区间内 $\Gamma$ 不变号，而相邻区间上 $\Gamma$ 的符号交替。
\end{enumerate}
\toHint{7.1}

\Q{7.2}{\lvII}{Legendre 倍元公式}
第~\bref{subsec:fe-two-symmetries}~节引用而未证明倍元公式
\[
  \Gamma(s)\,\Gamma\Bigl(s+\frac12\Bigr)=2^{1-2s}\sqrt\pi\,\Gamma(2s).
\]
用附录~\bref{app:weierstrass} 推论~\bref{cor:euler-form} 的 Gauss 极限形式
$\Gamma(z)=\lim_{n\to\infty}\dfrac{n!\,n^z}{z(z+1)\cdots(z+n)}$
证明它。（仓库「学习笔记」中《Γ 倍元公式的两种证明》给出另外两种证法。）
\toHint{7.2}

\Q{7.3}{\lvII}{$\chi(s)$ 在竖直方向的大小}
设 $\chi(s)=2^s\pi^{s-1}\sin(\pi s/2)\,\Gamma(1-s)$（式~(\bref{eq:chi_factor})）。用推论~\bref{cor:stirling-vertical} 证明：$\sigma$ 固定、$|t|\to\infty$ 时
\[
  |\chi(\sigma+it)|=\Bigl(\frac{|t|}{2\pi}\Bigr)^{1/2-\sigma}\bigl(1+O(1/|t|)\bigr).
\]
特别地，$|\chi(\tfrac12+it)|\to1$。（习题 9.5 将证明 $|\chi(\tfrac12+it)|=1$ 精确成立。）这里提前使用第~\bref{ch:riemann_zeta_continuation}~章的记号 $\chi$，题中已给出完整定义，只需第 7 章的工具；它是后文函数方程因子的预备计算。
\toHint{7.3}

\Q{7.4}{\lvII}{Stirling 公式的数值检验}
（(a) 为基础计算；(b)(c) 涉及复参数的余元公式与误差解释。）
\begin{enumerate}
  \item 对 $n=1,5,10,20$，计算 $n!\big/\bigl(\sqrt{2\pi n}\,(n/e)^n\bigr)$，并与 $1+\frac1{12n}$ 比较。
  \item 由余元公式证明：对实数 $t\ne0$，
  \[
    \bigl|\Gamma(\tfrac12+it)\bigr|^2=\frac{\pi}{\cosh\pi t},\qquad
    \bigl|\Gamma(it)\bigr|^2=\frac{\pi}{t\sinh\pi t}.
  \]
  \item 在 $t=1$ 与 $t=10$ 处，把 (b) 的精确值与推论~\bref{cor:stirling-vertical} 的主项 $2\pi|t|^{2\sigma-1}e^{-\pi|t|}$ 比较。误差为什么远小于推论所说的 $O(1/|t|)$？
\end{enumerate}
\toHint{7.4}

% ------------------------------------------------------------
\section*{第 \bref{ch:analytic_continuation_general} 章\quad 解析延拓的定义与唯一性}
\addcontentsline{toc}{section}{\texorpdfstring{第 \protect\bref{ch:analytic_continuation_general} 章\quad }{}解析延拓的定义与唯一性}

\Q{8.1}{\lvII}{收敛边界之外的部分和}
\begin{enumerate}
  \item 利用第一部习题 4.2 的表示 $\zeta(s)=\frac{s}{s-1}-s\int_1^\infty\{x\}x^{-s-1}\dd x$ 与有限 Abel 求和，证明对 $\Re s>0$、$s\ne1$ 及正整数 $N$，
  \[
    \sum_{n\le N}n^{-s}=\zeta(s)+\frac{N^{1-s}}{1-s}+s\int_N^\infty\{x\}x^{-s-1}\dd x,
  \]
  且最后一项的模不超过 $|s|N^{-\sigma}/\sigma$。
  \item 由此说明：$0<\sigma<1$ 时部分和的模满足 $\bigl|\sum_{n\le N}n^{-s}\bigr|\sim N^{1-\sigma}/|1-s|$；$\sigma=1$、$t\ne0$ 时部分和有界但不收敛（找出两条有不同极限的子列）。在这些点上失效的是级数，而不是 $\zeta$。
  \item 取 $s=\rho_1$（第一个非平凡零点，$\rho_1\approx\tfrac12+14.134725i$；计算时用高精度零点，如 \texttt{mpmath.zetazero(1)}，而不是这 $6$ 位小数），对 $N=100,1000,10000$ 计算 $\sum_{n\le N}n^{-s}$ 与 $\sum_{n\le N}n^{-s}-\frac{N^{1-s}}{1-s}$ 的模。
\end{enumerate}
\toHint{8.1}

\Q{8.2}{\lvI}{极点不产生多值性}
$\zeta$ 延拓到的 $\CC\setminus\{1\}$ 不是单连通区域。可以使用第~\bref{ch:riemann_zeta_continuation}~章的结论作为前提：积分构造给出了 $\CC\setminus\{1\}$ 上的单值亚纯函数，并在 $\Re s>1$ 上与原来的 $\zeta$ 一致。
\begin{enumerate}
  \item 说明为什么沿绕 $s=1$ 的不同路径作延拓，不会像 $\log z$ 那样得到不同的值；$s=1$ 附近的 Laurent 展开在这里说明什么，不说明什么？
  \item $\log\zeta(s)$ 能否在 $\CC\setminus\{1\}$ 上取单值解析分支？若把 $\zeta$ 的零点也全部除去呢？
\end{enumerate}
\toHint{8.2}

\Q{8.3}{\lvII}{另一个自然边界}
仿照第~\bref{sec:natural_boundary}~节的例子，证明 $f(z)=\sum_{n\ge0}z^{2^n}$ 以 $|z|=1$ 为自然边界。$f$ 满足函数关系 $f(z^2)=f(z)-z$；为什么这一关系不能像 $\Gamma(s)=\Gamma(s+1)/s$ 那样用来延拓（对照第~\bref{sec:fe_continuation_schema}~节）？
\toHint{8.3}

% ------------------------------------------------------------
\section*{第 \bref{ch:riemann_zeta_continuation} 章\quad $\zeta$ 函数的解析延拓与函数方程}
\addcontentsline{toc}{section}{\texorpdfstring{第 \protect\bref{ch:riemann_zeta_continuation} 章\quad }{}\texorpdfstring{$\zeta$}{ζ} 函数的解析延拓与函数方程}

\Q{9.1}{\lvI}{由 $\eta$ 计算 $\zeta(\tfrac12)$}
\begin{enumerate}
  \item 计算 $\eta(\tfrac12)$ 的部分和 $P_N=\sum_{n\le N}(-1)^{n-1}n^{-1/2}$（$N=10,100,1000$），再计算相邻部分和的平均 $\tfrac12(P_N+P_{N+1})$。哪一种收敛快？为什么？
  \item 由式~(\bref{eq:zeta_eta_continuation}) 求 $\zeta(\tfrac12)$。
\end{enumerate}
\toHint{9.1}

\Q{9.2}{\lvI}{负整数处的值}
由函数方程~(\bref{eq:asymmetric_fe_alt}) 与 $\zeta(2)=\pi^2/6$、$\zeta(4)=\pi^4/90$，求 $\zeta(-1)$、$\zeta(-3)$，并说明 $\zeta(-2)=\zeta(-4)=0$ 从何而来。
\toHint{9.2}

\Q{9.3}{\lvII}{Hankel 积分与 Bernoulli 数}
记 Bernoulli 数 $B_k$ 为展开 $\dfrac{z}{e^z-1}=\sum_{k\ge0}B_k\dfrac{z^k}{k!}$ 的系数（$B_0=1$，$B_1=-\tfrac12$，$B_2=\tfrac16$，$B_3=0$，$B_4=-\tfrac1{30}$）。
\begin{enumerate}
  \item 对非负整数 $m$，说明式~(\bref{eq:I_s}) 中 $I(-m)$ 的上下沿积分相互抵消，并求出 $I(-m)$。
  \item 由式~(\bref{eq:zeta_contour_gamma}) 证明 $\zeta(-m)=(-1)^m\dfrac{B_{m+1}}{m+1}$。
  \item 由此重新得到 $\zeta(0)$、$\zeta(-1)$、$\zeta(-3)$ 以及平凡零点。（这与第一部习题 3.3 是同一机制。）
\end{enumerate}
\toHint{9.3}

\Q{9.4}{\lvII}{为什么是 3--4--1}
定理~\bref{thm:no-zero-on-one} 的证明用到 $3+4\cos\theta+\cos2\theta\ge0$。以下固定实数 $t\ne0$。
\begin{enumerate}
  \item 若只用 $1+\cos\theta\ge0$，同样的论证给出什么不等式？能否排除 $\zeta(1+it)=0$？
  \item 解释 3--4--1 论证中 $|\zeta(\sigma+2it)|$ 一项为什么不会破坏结论。
  \item 取 $\sigma=1.05$、$1.01$ 与 $t=1$、$14.134725$，数值验证式~(\bref{eq:pnt-341})。
\end{enumerate}
\toHint{9.4}

\Q{9.5}{\lvII}{$\chi(s)\chi(1-s)=1$}
\begin{enumerate}
  \item 只用余元公式，由定义~(\bref{eq:chi_factor}) 直接证明 $\chi(s)\chi(1-s)=1$。这与把非对称函数方程连用两次有什么关系？
  \item 证明 $\chi(\bar s)=\overline{\chi(s)}$，并推出 $|\chi(\tfrac12+it)|=1$ 对一切实数 $t$ 精确成立。
  \item 由 (b) 说明：在临界线上 $|\zeta(\tfrac12+it)|=|\zeta(\tfrac12-it)|$，且 $\zeta(\tfrac12+it)$ 与 $\zeta(\tfrac12-it)$ 只差一个模为 $1$ 的因子。第~\bref{ch:riemann_numerical_approximation}~章的 Hardy $Z$ 函数正是由此构造的。
\end{enumerate}
\toHint{9.5}

\Q{9.6}{\lvIII}{$\zeta'(0)=-\tfrac12\log(2\pi)$}
利用函数方程~(\bref{eq:asymmetric_fe_alt})、$\zeta(1+u)=\frac1u+\gamma_E+O(u)$（$u\to0$；由第一部习题 4.2 的解析表示及其余项进一步展开得到），以及 $\Gamma'(1)=-\gamma_E$（第~\bref{subsec:gamma-poles}~节），求 $\zeta(s)$ 在 $s=0$ 处的 Taylor 展开到一次项，证明
\[
  \zeta'(0)=-\frac12\log(2\pi).
\]
\toHint{9.6}

\Q{9.7}{\lvI}{Riemann 第二法中的对称与极点分离}
记 $\omega(x)=\sum_{n\ge1}e^{-\pi n^2x}$，$H(s)=\int_1^\infty\bigl(x^{s/2-1}+x^{-(s+1)/2}\bigr)\omega(x)\dd x$。定理~\bref{thm:zeta-theta-integral} 给出
\[
  \pi^{-s/2}\Gamma\Bigl(\frac s2\Bigr)\zeta(s)=\frac1{s(s-1)}+H(s).
\]
\begin{enumerate}
  \item 不重复正文的证明，直接从 $H$ 的积分核读出 $H(s)=H(1-s)$，并说明 $H$ 在实轴上与临界线上都取实值。
  \item 数值计算 $H(2)$ 与 $H(\tfrac12)$，由上式核对 $\zeta(2)=\pi^2/6$ 与 $\zeta(\tfrac12)\approx-1.4603545$。再只用 $H(-1)=H(2)$ 求出 $\zeta(-1)$。
  \item 右端在 $s=0$ 与 $s=1$ 都有极点，但 $\zeta$ 只在 $s=1$ 有极点。结合 $1/\Gamma(s/2)$ 的零点说明原因；并说明 $\zeta(-2n)=0$（$n\ge1$）如何由同一式读出。
\end{enumerate}
\toHint{9.7}

% ============================================================
\newpage
\part*{二、提示}
\addcontentsline{toc}{part}{二、提示}
% ============================================================

\subsection*{第 7 章}
\Hint{7.1}(a) $\Gamma(s)=\Gamma(s+1)/s$。(c) 留数公式见定理~\bref{thm:gamma-meromorphic}；$\Gamma$ 在实轴上没有零点，故每个区间内不变号；跨过单极点时比较两侧。

\Hint{7.2}先证 $\prod_{k=0}^n(s+k)(s+k+\tfrac12)=2^{-2n-2}\prod_{k=0}^{2n+1}(2s+k)$；对 $\Gamma(2s)$ 取 Gauss 极限时用 $2n+1$ 项。比值的极限与 $s$ 无关，再令 $s=\tfrac12$ 定常数。

\Hint{7.3}$|\sin(a+ib)|^2=\sin^2a+\sinh^2b$；$\Gamma(1-s)$ 的实部为 $1-\sigma$。指数因子 $e^{\pm\pi|t|/2}$ 恰好相消。

\Hint{7.4}(b) $\Gamma(\bar s)=\overline{\Gamma(s)}$；$\sin(\pi(\tfrac12+it))=\cosh\pi t$，$\sin(i\pi t)=i\sinh\pi t$；$\Gamma(1-it)=-it\,\Gamma(-it)$。

\subsection*{第 8 章}
\Hint{8.1}(a) 有限 Abel 求和：$\sum_{n\le N}n^{-s}=N^{1-s}+s\int_1^N\lfloor x\rfloor x^{-s-1}\dd x$。(c) 在零点处 $\zeta(\rho_1)=0$。

\Hint{8.2}(a) 第 8 章正文已指出：唯一性定理本身不保证沿不同区域链的延拓能拼成整体单值函数；这里的整体相容性来自第 9 章的整体构造。(b) 若 $e^L=\zeta$，$\zeta$ 能取零值吗？绕 $s=1$ 的小圆上，$\frac1{2\pi i}\oint\zeta'/\zeta\,\dd s$ 等于多少？

\Hint{8.3}在 $2^k$ 次单位根 $\omega$ 处，$n\ge k$ 时 $\omega^{2^n}=1$。映射 $z\mapsto z^2$ 把单位圆盘映入自身。

\subsection*{第 9 章}
\Hint{9.1}交错级数的部分和在极限两侧来回摆动，相邻两项的平均抵消了主要的摆动。要得到误差阶，把 $\eta(\tfrac12)-\frac12(P_N+P_{N+1})$ 写成相邻项之差 $a_n-a_{n+1}$（$a_n=n^{-1/2}$）组成的交错级数。

\Hint{9.2}$\sin(-\pi/2)=-1$，$\sin(-3\pi/2)=1$，$\Gamma(2)=1$，$\Gamma(4)=6$。

\Hint{9.3}(a) 整数 $s$ 时 $(-z)^{s-1}$ 单值。$\dfrac1{e^z-1}=\sum_{k\ge0}B_k\dfrac{z^{k-1}}{k!}$。(b) $\Gamma(1+m)=m!$。

\Hint{9.4}(a) 若 $1+it$ 是 $m$ 阶零点，$\zeta(\sigma)|\zeta(\sigma+it)|\asymp(\sigma-1)^{m-1}$。

\Hint{9.5}(a) $\sin\frac{\pi(1-s)}2=\cos\frac{\pi s}2$，$2\sin\frac{\pi s}2\cos\frac{\pi s}2=\sin\pi s$。(b) 在 $s=\tfrac12+it$ 处，$1-s=\bar s$。

\Hint{9.6}$\chi(s)=\frac s2\bigl(1+s(\log2\pi+\gamma_E)\bigr)+O(s^3)$；$\zeta(1-s)=-\frac1s+\gamma_E+O(s)$。

\Hint{9.7}(a) $s\mapsto1-s$ 使两个幂函数 $x^{s/2-1}$ 与 $x^{-(s+1)/2}$ 互换。(b) $\Gamma(1)=1$，$\Gamma(-\frac12)=-2\sqrt\pi$。(c) $1/\Gamma(s/2)$ 在 $s=0,-2,-4,\dots$ 有单零点，而右端在 $s=-2n$（$n\ge1$）处有限。

% ============================================================
\newpage
\part*{三、解答}
\addcontentsline{toc}{part}{三、解答}
% ============================================================

\subsection*{第 7 章}

\Sol{7.1}{负半整数处的值与留数}
(a) $\Gamma(\tfrac52)=\tfrac32\cdot\tfrac12\Gamma(\tfrac12)=\tfrac34\sqrt\pi\approx1.32934$；
$\Gamma(-\tfrac12)=\Gamma(\tfrac12)/(-\tfrac12)=-2\sqrt\pi\approx-3.54491$；
$\Gamma(-\tfrac32)=\Gamma(-\tfrac12)/(-\tfrac32)=\tfrac43\sqrt\pi\approx2.36327$。

(b) $\Gamma(-\tfrac12)\Gamma(\tfrac32)=-2\sqrt\pi\cdot\tfrac{\sqrt\pi}2=-\pi$，而 $\pi/\sin(-\pi/2)=-\pi$，一致。

(c) $\Res_{s=-1}\Gamma=-1$，$\Res_{s=-2}\Gamma=\tfrac12$。在实轴上，$\Gamma$ 在 $(-n-1,-n)$ 内连续、取实值且没有零点（定理~\bref{thm:gamma-no-zeros}），故符号不变；而在单极点 $-n$ 两侧，$\Gamma(s)\approx\frac{(-1)^n/n!}{s+n}$ 的符号相反。所以 $\Gamma$ 在相邻区间 $(-n-1,-n)$ 与 $(-n,-n+1)$ 上符号相反：在 $(-1,0)$ 上为负，$(-2,-1)$ 上为正，依此交替。(a) 中 $\Gamma(-\tfrac12)<0$、$\Gamma(-\tfrac32)>0$ 与此一致。

\Sol{7.2}{Legendre 倍元公式}
记 $P_n(z)=z(z+1)\cdots(z+n)$。由于 $(s+k)(s+k+\tfrac12)=\tfrac14(2s+2k)(2s+2k+1)$，
\[
  P_n(s)\,P_n\bigl(s+\tfrac12\bigr)=2^{-2n-2}\prod_{k=0}^{2n+1}(2s+k)=2^{-2n-2}P_{2n+1}(2s).
\]
设 $2s\notin\{0,-1,-2,\dots\}$。对 $\Gamma(s)$、$\Gamma(s+\tfrac12)$ 取 $n$ 项的 Gauss 极限，对 $\Gamma(2s)$ 取 $m=2n+1$ 项的 Gauss 极限：
\[
\begin{aligned}
  \frac{\Gamma(s)\Gamma(s+\frac12)}{\Gamma(2s)}
  &=\lim_{n\to\infty}\frac{n!^2\,n^{2s+1/2}\,2^{2n+2}\big/P_{2n+1}(2s)}{(2n+1)!\,(2n+1)^{2s}\big/P_{2n+1}(2s)}\\
  &=\lim_{n\to\infty}\Bigl(\frac{n}{2n+1}\Bigr)^{2s}\cdot\frac{n!^2\,2^{2n+2}\sqrt n}{(2n+1)!}.
\end{aligned}
\]
第一个因子趋于 $2^{-2s}$。左端的极限存在，所以第二个因子（与 $s$ 无关）的极限 $C$ 也存在，且
$\Gamma(s)\Gamma(s+\tfrac12)=2^{-2s}C\,\Gamma(2s)$。令 $s=\tfrac12$：$\Gamma(\tfrac12)\Gamma(1)=\tfrac12C\,\Gamma(1)$，得 $C=2\sqrt\pi$。于是
$\Gamma(s)\Gamma(s+\tfrac12)=2^{1-2s}\sqrt\pi\,\Gamma(2s)$。两端是亚纯函数，由恒等定理对一切 $s$ 成立（在极点处按亚纯恒等式理解）。

\Sol{7.3}{$\chi(s)$ 在竖直方向的大小}
记 $s=\sigma+it$。由 $|\sin(a+ib)|^2=\sin^2a+\sinh^2b$，
\[
  \Bigl|\sin\frac{\pi s}2\Bigr|^2=\sin^2\frac{\pi\sigma}2+\sinh^2\frac{\pi t}2=\frac14e^{\pi|t|}\bigl(1+O(e^{-\pi|t|})\bigr),
\]
故 $|\sin(\pi s/2)|=\tfrac12e^{\pi|t|/2}\bigl(1+O(e^{-\pi|t|})\bigr)$。$1-s=(1-\sigma)-it$，由推论~\bref{cor:stirling-vertical}（实部固定为 $1-\sigma$），
\[
  |\Gamma(1-s)|=\sqrt{2\pi}\,|t|^{1/2-\sigma}e^{-\pi|t|/2}\bigl(1+O(1/|t|)\bigr).
\]
又 $|2^s\pi^{s-1}|=2^\sigma\pi^{\sigma-1}$。相乘，指数因子相消：
\[
  |\chi(s)|=2^\sigma\pi^{\sigma-1}\cdot\tfrac12\cdot\sqrt{2\pi}\,|t|^{1/2-\sigma}\bigl(1+O(1/|t|)\bigr)
  =(2\pi)^{\sigma-1/2}|t|^{1/2-\sigma}\bigl(1+O(1/|t|)\bigr),
\]
即所求。$\sigma=\tfrac12$ 时主项为 $1$。

\Sol{7.4}{Stirling 公式的数值检验}
(a)
\begin{center}
\begin{tabular}{rrrr}
\toprule
$n$ & $n!$ & $n!/\bigl(\sqrt{2\pi n}(n/e)^n\bigr)$ & $1+\frac1{12n}$\\
\midrule
$1$ & $1$ & $1.084438$ & $1.083333$\\
$5$ & $120$ & $1.016784$ & $1.016667$\\
$10$ & $3628800$ & $1.008365$ & $1.008333$\\
$20$ & $2.4329\times10^{18}$ & $1.004175$ & $1.004167$\\
\bottomrule
\end{tabular}
\end{center}
比值趋于 $1$，相对误差约为 $\frac1{12n}$。这正是 $E_\Gamma$ 的首项：可以证明 $E_\Gamma(s)=\frac1{12s}+O(|s|^{-3})$，定理~\bref{thm:stirling} 只给出它的上界 $O(1/|s|)$。

(b) $\Gamma(\tfrac12-it)=\overline{\Gamma(\tfrac12+it)}$，由余元公式
$|\Gamma(\tfrac12+it)|^2=\Gamma(\tfrac12+it)\Gamma(\tfrac12-it)=\pi/\sin\bigl(\pi(\tfrac12+it)\bigr)=\pi/\cosh\pi t$。
又 $\Gamma(1-it)=-it\,\Gamma(-it)$，故
$|\Gamma(it)|^2=\Gamma(it)\Gamma(-it)=\dfrac{\Gamma(it)\Gamma(1-it)}{-it}=\dfrac{\pi}{-it\sin(i\pi t)}=\dfrac{\pi}{-it\cdot i\sinh\pi t}=\dfrac{\pi}{t\sinh\pi t}$。

(c) $t=1$：$|\Gamma(\tfrac12+i)|^2\approx0.271015$，主项 $2\pi e^{-\pi}\approx0.271521$；$|\Gamma(i)|^2\approx0.272029$，主项同为 $0.271521$。相对误差约 $0.2\%$。$t=10$：两个精确值分别约为 $1.42697\times10^{-13}$ 与 $1.42697\times10^{-14}$，与各自的主项的相对差约为 $5.2\times10^{-28}$。

原因可直接由 (b) 的精确式读出。令 $q=e^{-2\pi|t|}$，则
\[
  \bigl|\Gamma(\tfrac12+it)\bigr|^2=\frac{2\pi e^{-\pi|t|}}{1+q},\qquad
  \bigl|\Gamma(it)\bigr|^2=\frac{2\pi|t|^{-1}e^{-\pi|t|}}{1-q},
\]
与主项之比分别是 $1/(1+q)$ 与 $1/(1-q)$，相对误差约为 $q=e^{-2\pi|t|}$，是指数级小量（$t=1$ 时 $q\approx0.0019$，$t=10$ 时 $q\approx5.2\times10^{-28}$）。推论中的 $O(1/|t|)$ 只是对一般 $\sigma$ 成立的上界，不说明实际误差的大小；在这两条特殊竖线上，余元公式把 $|\Gamma|^2$ 化成了初等函数，误差的结构因而一目了然。

\subsection*{第 8 章}

\Sol{8.1}{收敛边界之外的部分和}
(a) 有限 Abel 求和（$A(x)=\lfloor x\rfloor$，$A(N)=N$）：
\[
  \sum_{n\le N}n^{-s}=N\cdot N^{-s}+s\int_1^N\lfloor x\rfloor x^{-s-1}\dd x
  =N^{1-s}+\frac{s}{1-s}\bigl(N^{1-s}-1\bigr)-s\int_1^N\{x\}x^{-s-1}\dd x .
\]
前两项之和为 $\frac{N^{1-s}}{1-s}+\frac{s}{s-1}$。再减去第一部习题 4.2 的
$\zeta(s)=\frac{s}{s-1}-s\int_1^\infty\{x\}x^{-s-1}\dd x$（对 $\Re s>0$、$s\ne1$ 成立），得
\[
  \sum_{n\le N}n^{-s}-\zeta(s)=\frac{N^{1-s}}{1-s}+s\int_N^\infty\{x\}x^{-s-1}\dd x .
\]
由 $0\le\{x\}<1$，最后一项的模不超过 $|s|\int_N^\infty x^{-\sigma-1}\dd x=|s|N^{-\sigma}/\sigma$。

(b) $0<\sigma<1$ 时 $|N^{1-s}/(1-s)|=N^{1-\sigma}/|1-s|\to\infty$，其余两项有界，故部分和的模满足 $\bigl|\sum_{n\le N}n^{-s}\bigr|\sim N^{1-\sigma}/|1-s|$。这是关于模的渐近式，并不说明实部、虚部或模对每个 $N$ 单调增长。

$\sigma=1$、$t\ne0$ 时，$\frac{N^{1-s}}{1-s}=\frac{N^{-it}}{-it}$ 的模恒为 $1/|t|$，余项趋于零，故部分和有界，并逐渐靠近以 $\zeta(1+it)$ 为圆心、半径 $1/|t|$ 的圆。取 $N_k=\lfloor e^{2\pi k/|t|}\rfloor$ 与 $M_k=\lfloor e^{(2\pi k+\pi)/|t|}\rfloor$，则 $t\log N_k$ 与 $2\pi k\operatorname{sgn}t$ 之差趋于 $0$，故 $N_k^{-it}\to1$；同理 $M_k^{-it}\to-1$。两条子列的部分和分别趋于 $\zeta(1+it)\pm\frac1{-it}$，极限不同，所以部分和不收敛。$\zeta(s)$ 本身在这些点有确定的值，只是级数不再表示它。

(c) $\rho_1$ 处 $\zeta(\rho_1)=0$：
\begin{center}
\begin{tabular}{rrr}
\toprule
$N$ & $\bigl|\sum_{n\le N}n^{-\rho_1}\bigr|$ & $\bigl|\sum_{n\le N}n^{-\rho_1}-\frac{N^{1-\rho_1}}{1-\rho_1}\bigr|$\\
\midrule
$100$ & $0.709$ & $0.0500$\\
$1000$ & $2.236$ & $0.0158$\\
$10000$ & $7.071$ & $0.0050$\\
\bottomrule
\end{tabular}
\end{center}
部分和本身按 $\sqrt N$ 增长（$\approx\sqrt N/|1-\rho_1|$），扣除显式的修正项后迅速趋于 $\zeta(\rho_1)=0$。剩余误差恰约为 $\tfrac12N^{-1/2}$，远小于 (a) 的粗略上界 $|s|N^{-1/2}/\sigma\approx28N^{-1/2}$。它的来源要用更高阶的 Euler--Maclaurin 展开说明：下一项为 $\tfrac12N^{-s}$，其模为 $\tfrac12N^{-1/2}$。这是补充解释，(a) 的粗界本身并不给出这个系数。

\Sol{8.2}{极点不产生多值性}
(a) 逻辑顺序是「整体构造保证全局相容，局部展开说明极点性质」，不能倒过来。

全局：第 9 章的积分构造直接给出 $\CC\setminus\{1\}$ 上的一个单值亚纯函数，它在 $\Re s>1$ 上等于原来的 $\zeta$。沿任何允许的路径作延拓，每一步得到的函数元素都与这个整体函数在重叠处一致（延拓唯一性），所以终点处的值就是整体函数在该点的值，与路径无关。仅凭唯一性定理不能得出这一点：第 8 章正文已指出，沿不同区域链的延拓未必能拼成整体单值函数，$\log z$ 就是例子。

局部：$s=1$ 是孤立奇点，$\zeta$ 在去心圆盘 $0<|s-1|<r$ 内有 Laurent 展开 $\frac1{s-1}+\sum_{k\ge0}c_k(s-1)^k$（定理~\bref{thm:laurent}），主部只有一项。这说明 $s=1$ 是单极点，而不是像 $\log z$ 在 $0$ 那样的对数型分支点：绕它一周，函数值不变。但它只描述 $s=1$ 附近，单独不能保证远处的路径相容。

(b) 不能，而且有两层障碍。

第一层：$\CC\setminus\{1\}$ 中含有 $\zeta$ 的零点。若存在解析函数 $L$ 使 $e^L=\zeta$，则 $\zeta=e^L$ 处处不为零，矛盾。所以在包含零点的区域上，$\log\zeta$ 根本无从定义。

第二层：即使把全部零点也除去，剩余区域上仍没有单值解析的对数。取绕 $s=1$ 的充分小正向圆周，其内部只有极点、没有零点，由辐角原理 $\frac1{2\pi i}\oint\frac{\zeta'}{\zeta}\dd s=-1$；若解析对数 $L$ 存在，则 $L'=\zeta'/\zeta$，这个积分应为 $0$。等价地说，$\log\zeta$ 在 $s=1$ 附近像 $-\log(s-1)$，绕一周改变 $-2\pi i$。

所以第 5 章只在无零点的单连通半平面 $\Re s>1$ 上定义 $\log\zeta$（备注~\bref{rem:log_zeta_branch}）。

\Sol{8.3}{另一个自然边界}
$|z|<1$ 时级数绝对收敛，$f$ 解析。取 $\omega=e^{2\pi ip/2^k}$。当 $n\ge k$ 时 $2^k\mid2^n$，$\omega^{2^n}=1$，故对 $0<r<1$，
\[
  f(r\omega)=\sum_{n<k}(r\omega)^{2^n}+\sum_{n\ge k}r^{2^n}.
\]
前一部分有界；后一部分各项为正，对任意 $M$，前 $M$ 项之和在 $r\to1^-$ 时趋于 $M$，故 $|f(r\omega)|\to\infty$。与第~\bref{sec:natural_boundary}~节同样的论证说明 $\omega$ 不能跨越；形如 $e^{2\pi ip/2^k}$ 的点在单位圆上稠密，可跨越的边界点构成开集而不含这些点，只能为空集。

函数关系：$f(z^2)=\sum_{n\ge0}z^{2^{n+1}}=f(z)-z$。用它延拓，需要把新区域中的点映回已知区域；但 $\Phi(z)=z^2$ 把单位圆盘映入自身，把圆外的点映到更远的圆外，从不把外面的点带回圆盘内。$\Gamma$ 的情形则不同：$\Phi(s)=s+1$ 把左边的竖带映入已经定义好的右边区域。函数关系能否用于延拓，取决于 $\Phi$ 能否把新点送回已知区域。

\subsection*{第 9 章}

\Sol{9.1}{由 $\eta$ 计算 $\zeta(\tfrac12)$}
(a)
\begin{center}
\begin{tabular}{rrr}
\toprule
$N$ & $P_N$ & $\tfrac12(P_N+P_{N+1})$\\
\midrule
$10$ & $0.45073$ & $0.601481$\\
$100$ & $0.55502$ & $0.604775$\\
$1000$ & $0.58909$ & $0.604895$\\
\bottomrule
\end{tabular}
\end{center}
真值 $\eta(\tfrac12)\approx0.6048986$。$P_N$ 的误差约为 $\pm\tfrac12N^{-1/2}$，随 $N$ 的奇偶在极限两侧交替。表中的 $N=10,100,1000$ 都是偶数，最后一项为负，所以表中的 $P_N$ 全在真值下方；「两侧交替」是对全部整数 $N$ 而言的。

相邻平均的误差可以严格估计。记 $a_n=n^{-1/2}$，$A_N=\frac12(P_N+P_{N+1})=P_N+\frac12(-1)^Na_{N+1}$，$b_j=a_{N+1+j}-a_{N+2+j}$。则
\[
  \eta(\tfrac12)-A_N=(-1)^N\Bigl(\tfrac12a_{N+1}-a_{N+2}+a_{N+3}-\cdots\Bigr)
  =\frac{(-1)^N}2\sum_{j\ge0}(-1)^jb_j .
\]
由于 $a_n$ 是凸函数在整数点的值，$b_j$ 为正、递减且趋于零，交错级数的余项不超过首项，故
\[
  \bigl|\eta(\tfrac12)-A_N\bigr|\le\frac{a_{N+1}-a_{N+2}}2=O(N^{-3/2}).
\]
（$a_{N+1}-a_{N+2}\approx\frac12N^{-3/2}$，所以误差约不超过 $\frac14N^{-3/2}$，与表中数据相符。）

(b) $\zeta(\tfrac12)=\eta(\tfrac12)/(1-2^{1/2})=0.6048986/(-0.4142136)\approx-1.4603545$。

\Sol{9.2}{负整数处的值}
由 $\zeta(s)=2^s\pi^{s-1}\sin(\pi s/2)\Gamma(1-s)\zeta(1-s)$：
\[
  \zeta(-1)=2^{-1}\pi^{-2}\cdot\sin(-\tfrac\pi2)\cdot\Gamma(2)\cdot\zeta(2)=\frac1{2\pi^2}\cdot(-1)\cdot1\cdot\frac{\pi^2}6=-\frac1{12},
\]
\[
  \zeta(-3)=2^{-3}\pi^{-4}\cdot\sin(-\tfrac{3\pi}2)\cdot\Gamma(4)\cdot\zeta(4)=\frac1{8\pi^4}\cdot1\cdot6\cdot\frac{\pi^4}{90}=\frac1{120}.
\]
$s=-2n$（$n\ge1$；$s=0$ 不在其列，$\zeta(0)=-\frac12$）时 $\sin(-n\pi)=0$，而 $\Gamma(1+2n)\zeta(1+2n)$ 有限，故 $\zeta(-2n)=0$：这就是平凡零点（第~\bref{subsec:trivial-zeros-simple}~节）。

\Sol{9.3}{Hankel 积分与 Bernoulli 数}
(a) $s=-m$ 时 $(-z)^{s-1}=(-z)^{-m-1}$ 是 $z$ 的单值函数，上下沿取值相同、方向相反，积分抵消；剩下绕原点一周的小圆积分（半径 $<2\pi$，圆内只有极点 $z=0$）。由
$\dfrac1{e^z-1}=\sum_{k\ge0}B_k\dfrac{z^{k-1}}{k!}$，
$(-z)^{-m-1}/(e^z-1)$ 的 $z^{-1}$ 系数为 $(-1)^{m+1}$ 乘以 $\frac1{e^z-1}$ 的 $z^m$ 系数，即 $(-1)^{m+1}\frac{B_{m+1}}{(m+1)!}$。故
\[
  I(-m)=2\pi i\,(-1)^{m+1}\frac{B_{m+1}}{(m+1)!}.
\]
$m=0$ 时为 $2\pi i\cdot(-1)\cdot(-\tfrac12)=\pi i$，与正文的 $I(0)$ 一致。

(b) $s=-m$ 处 $\Gamma(1-s)=m!$ 有限，式~(\bref{eq:zeta_contour_gamma}) 直接适用：
\[
  \zeta(-m)=-\frac{m!}{2\pi i}\,I(-m)=(-1)^m\frac{m!\,B_{m+1}}{(m+1)!}=(-1)^m\frac{B_{m+1}}{m+1}.
\]

(c) $\zeta(0)=B_1=-\tfrac12$；$\zeta(-1)=-B_2/2=-\tfrac1{12}$；$\zeta(-3)=-B_4/4=\tfrac1{120}$，与习题 9.2 一致。$m=2n\ge2$ 时 $B_{2n+1}=0$（因为 $\frac z{e^z-1}+\frac z2$ 是偶函数，大于 $1$ 的奇数下标 Bernoulli 数都为零），故 $\zeta(-2n)=0$。

\Sol{9.4}{为什么是 3--4--1}
(a) 由 $\log|\zeta(\sigma+it)|=\sum\frac{\Lambda(n)}{n^\sigma\log n}\cos(t\log n)$ 与 $1+\cos\theta\ge0$，得 $\zeta(\sigma)|\zeta(\sigma+it)|\ge1$（$\sigma>1$）。若 $1+it$ 是 $m$ 阶零点，左端 $\asymp(\sigma-1)^{-1}(\sigma-1)^m=(\sigma-1)^{m-1}$。$m=1$ 时左端有界，不产生矛盾；只有 $m\ge2$ 才矛盾。所以这一不等式只能排除重零点，排除不了单零点。

(b) $|\zeta(\sigma+2it)|$ 出现在左端的乘积中，它可能很大也可能很小。论证只需要它在 $\sigma\to1^+$ 时\textbf{有界}：因为 $t\ne0$，$1+2it\ne1$，$\zeta$ 在其附近解析，故有界。于是左端 $\ll(\sigma-1)^{-3}(\sigma-1)^{4m}\le\sigma-1\to0$，与 $\ge1$ 矛盾。三角式中各系数 $3,4,1$ 的作用正在于：$\cos\theta$ 的系数 $4$ 大于 $1$ 的系数 $3$，使零点的阶 $4m$ 压过极点的阶 $3$。

(c)
\begin{center}
\begin{tabular}{rrrr}
\toprule
$\sigma$ & $t$ & $\zeta(\sigma)^3|\zeta(\sigma+it)|^4|\zeta(\sigma+2it)|$ & $\zeta(\sigma)|\zeta(\sigma+it)|$\\
\midrule
$1.05$ & $1$ & $9797$ & $23.10$\\
$1.05$ & $14.134725$ & $254.4$ & $7.248$\\
$1.01$ & $1$ & $1.039\times10^6$ & $110.7$\\
$1.01$ & $14.134725$ & $2.381\times10^4$ & $33.34$\\
\bottomrule
\end{tabular}
\end{center}
各值均大于 $1$。注意 $t=14.134725$ 是第一个零点的纵坐标，但零点在 $\sigma=\tfrac12$ 上，离直线 $\sigma=1$ 很远，所以乘积并不小。

\Sol{9.5}{$\chi(s)\chi(1-s)=1$}
(a) $\chi(1-s)=2^{1-s}\pi^{-s}\sin\frac{\pi(1-s)}2\Gamma(s)=2^{1-s}\pi^{-s}\cos\frac{\pi s}2\Gamma(s)$。于是
\[
  \chi(s)\chi(1-s)=2\pi^{-1}\sin\frac{\pi s}2\cos\frac{\pi s}2\,\Gamma(s)\Gamma(1-s)=\frac{\sin\pi s}{\pi}\cdot\frac{\pi}{\sin\pi s}=1
\]
（先在 $s\notin\mathbb Z$ 处计算，再按亚纯恒等式推广）。非对称函数方程连用两次：$\zeta(s)=\chi(s)\zeta(1-s)=\chi(s)\chi(1-s)\zeta(s)$，这一恒等式正说明两次反射回到原处，彼此相容。

(b) $\chi$ 的定义中只有实系数的函数（$2^s$、$\pi^{s-1}$、$\sin$、$\Gamma$ 在实轴上取实值），故 $\chi(\bar s)=\overline{\chi(s)}$。在 $s=\tfrac12+it$ 处 $1-s=\tfrac12-it=\bar s$，所以
$1=\chi(s)\chi(\bar s)=|\chi(s)|^2$。这比习题 7.3 的渐近结论更强：在临界线上 $|\chi|=1$ 对一切 $t$ 精确成立。

(c) 由 $\zeta(s)=\chi(s)\zeta(1-s)$ 与 $1-s=\bar s$，$\zeta(\tfrac12+it)=\chi(\tfrac12+it)\,\zeta(\tfrac12-it)$，而 $|\chi(\tfrac12+it)|=1$。又 $\zeta(\tfrac12-it)=\overline{\zeta(\tfrac12+it)}$，所以 $\zeta(\tfrac12+it)$ 与其共轭只差一个模为 $1$ 的因子。函数 $t\mapsto\chi(\tfrac12+it)$ 在整条实轴上连续、不为零且模为 $1$；实轴单连通，所以可以选取连续的实函数 $\vartheta$，使 $\chi(\tfrac12+it)=e^{-2i\vartheta(t)}$。这样的 $\vartheta$ 只确定到相差 $\pi$ 的整数倍，需要一个归一化：由 $\zeta(\tfrac12)=\chi(\tfrac12)\zeta(\tfrac12)$ 与 $\zeta(\tfrac12)\ne0$ 得 $\chi(\tfrac12)=1$，取 $\vartheta(0)=0$，这与第~\bref{ch:riemann_numerical_approximation}~章的相位约定一致；否则所构造的 $Z$ 会有一个整体正负号的不确定性。（逐点各取一个辐角，并不能自动拼成一条连续相位。）令 $Z(t)=e^{i\vartheta(t)}\zeta(\tfrac12+it)$，则
\[
  \overline{Z(t)}=e^{-i\vartheta(t)}\zeta(\tfrac12-it)=e^{-i\vartheta(t)}\,\overline{\chi(\tfrac12+it)}\,\zeta(\tfrac12+it)=e^{-i\vartheta(t)}e^{2i\vartheta(t)}\zeta(\tfrac12+it)=Z(t),
\]
其中第二个等号用了 $\zeta(\tfrac12-it)=\chi(\tfrac12-it)\zeta(\tfrac12+it)$ 与 $\chi(\tfrac12-it)=\overline{\chi(\tfrac12+it)}$。所以 $Z(t)$ 取实值且 $|Z(t)|=|\zeta(\tfrac12+it)|$，这就是 Hardy $Z$ 函数的构造。以上推导始终使用乘法形式 $\zeta(s)=\chi(s)\zeta(1-s)$，没有写成商 $\zeta(s)/\zeta(\bar s)$，所以在 $\zeta$ 的零点处同样成立（那里商是 $0/0$，无定义）。

\Sol{9.6}{$\zeta'(0)=-\tfrac12\log(2\pi)$}
在 $s=0$ 附近展开 $\chi(s)=2^s\pi^{s-1}\sin(\pi s/2)\Gamma(1-s)$ 的各因子：
\[
\begin{gathered}
  2^s\pi^{s-1}=\frac1\pi\bigl(1+s\log(2\pi)+O(s^2)\bigr),\qquad
  \sin\frac{\pi s}2=\frac{\pi s}2+O(s^3),\\
  \Gamma(1-s)=1-\Gamma'(1)s+O(s^2)=1+\gamma_E s+O(s^2).
\end{gathered}
\]
故 $\chi(s)=\frac s2\bigl(1+s(\log2\pi+\gamma_E)\bigr)+O(s^3)$。由 $\zeta(1+u)=\frac1u+\gamma_E+O(u)$（取 $u=-s$），$\zeta(1-s)=-\frac1s+\gamma_E+O(s)$。相乘：
\[
  \zeta(s)=\chi(s)\zeta(1-s)=\frac s2\Bigl(-\frac1s+\gamma_E\Bigr)+\frac{s^2}2(\log2\pi+\gamma_E)\Bigl(-\frac1s\Bigr)+O(s^2)
  =-\frac12-\frac s2\log(2\pi)+O(s^2).
\]
常数项 $-\tfrac12$ 与式~(\bref{eq:zeta-at-zero}) 一致，一次项系数即 $\zeta'(0)=-\tfrac12\log(2\pi)\approx-0.918939$（数值核对一致）。$\gamma_E$ 在两处出现、恰好相消。

\Sol{9.7}{Riemann 第二法中的对称与极点分离}
(a) 把 $s$ 换成 $1-s$：$x^{(1-s)/2-1}=x^{-(s+1)/2}$，$x^{-((1-s)+1)/2}=x^{s/2-1}$，两个幂函数恰好互换，其和不变，故 $H(1-s)=H(s)$。对称函数方程~(\bref{eq:symmetric_fe}) 就是由此直接读出的（极点项 $\frac1{s(s-1)}$ 同样在 $s\mapsto1-s$ 下不变）。
$s$ 为实数时被积函数为实，$H(s)\in\RR$；$s=\frac12+it$ 时 $x^{s/2-1}=x^{-3/4}x^{it/2}$，$x^{-(s+1)/2}=x^{-3/4}x^{-it/2}$，和为 $2x^{-3/4}\cos(\frac t2\log x)$，也是实数。

(b) 数值积分得 $H(2)\approx0.0235988$，$H(\tfrac12)\approx0.0230338$。
$s=2$：左端 $\pi^{-1}\Gamma(1)\zeta(2)$，右端 $\frac12+H(2)\approx0.5235988=\frac\pi6$，故 $\zeta(2)=\frac{\pi^2}6$。
$s=\frac12$：右端 $-4+H(\tfrac12)\approx-3.9769662$；除以 $\pi^{-1/4}\Gamma(\tfrac14)\approx2.7232893$，得 $\zeta(\tfrac12)\approx-1.4603545$。
$s=-1$：$H(-1)=H(2)$，右端 $\frac12+H(2)=\frac\pi6$；左端 $\pi^{1/2}\Gamma(-\tfrac12)\zeta(-1)=-2\pi\zeta(-1)$，故 $\zeta(-1)=-\frac1{12}$，与习题 9.2 一致。

(c) 右端在 $s=0$、$1$ 的极点属于完成函数 $\pi^{-s/2}\Gamma(s/2)\zeta(s)$，不都是 $\zeta$ 的极点。改写为 $\zeta(s)=\pi^{s/2}\frac1{\Gamma(s/2)}\bigl(\frac1{s(s-1)}+H(s)\bigr)$：$1/\Gamma(s/2)$ 在 $s=0$ 有单零点，恰好抵消 $\frac1{s(s-1)}$ 在 $0$ 的单极点，所以 $\zeta$ 在 $0$ 解析（并得 $\zeta(0)=-\frac12$）；在 $s=1$，$1/\Gamma(\frac12)=1/\sqrt\pi\ne0$，单极点保留。
$s=-2n$（$n\ge1$）时右端 $\frac1{s(s-1)}+H(s)$ 有限，而 $1/\Gamma(s/2)$ 在此为零，故 $\zeta(-2n)=0$：平凡零点来自 $\Gamma(s/2)$ 的极点。

\end{document}
